\section{Experimental Setup}
\label{sec:setup}

\subsection{Experimental Questions}

\textbf{Question 1 (h-e1).} Does ratio reward produce detectably different GRPO advantage variance than binary reward in groups representative of early-training conditions?

\textbf{Question 2 (h-m1).} Does the gradient signal difference (Q1) translate to measurable policy-level performance improvements after 1,000 GRPO training steps?

\subsection{Evaluation Metrics}

\textbf{Advantage variance.} $\text{Var}(\{r_i - \text{mean}(r)\})$, directly measuring the group-level gradient signal.

\textbf{FractionPartialCallback.} Fraction of completions per batch with $0 < k < n$ test cases passing. If \texttt{fraction\_partial = 0}, ratio and binary rewards are equivalent.

\textbf{Clipped ratio.} Fraction of completions reaching \texttt{max\_new\_tokens} limit. \texttt{clipped\_ratio = 1.0} indicates all completions are truncated, precluding test execution.

\textbf{HumanEval pass@1.} Greedy-decode pass@1 on 164 problems \citep{chen2021humaneval}, evaluated at training checkpoints.

\subsection{h-e1 Components}

\begin{enumerate}
    \item \textbf{Synthetic group computation.} Direct computation of advantages for group $[0,0,1,2,0,3,0,0]$ (8 completions, 5 test cases) under both reward functions.
    \item \textbf{Scale simulation.} 1,000 GRPO groups ($G=8$, $n=5$, $p_{\text{pass}}=0.1$ per test); record fraction with ratio$\neq$binary signal.
    \item \textbf{Smoke test.} 3 actual GRPO training steps; verify exit code 0, gradient norm logging, checkpoint saving.
\end{enumerate}

\subsection{h-m1 Design}

Full GRPO training for 1,000 steps under both conditions simultaneously on separate GPUs (binary: GPU 3; ratio: GPU 4). Checkpoints at steps 200, 400, 600, 800, 1,000. HumanEval evaluated at all checkpoints. FractionPartialCallback fires at every step as the primary diagnostic.
